% !PRINTED 245
\weberpage{245}

\weq{3}{a_1 < 0, \quad a_2 > 0, \quad a_3 < 0;}
these conditions are necessary in order that all three roots be
positive. But they are also sufficient; for if one or three roots
are negative, then $a_3 > 0$. If, however, two roots, say
$\beta, \gamma$, are negative, then either
\[
a_1 \geqq 0 \quad \text{or} \quad \alpha > -(\beta + \gamma).
\]
But in the latter case it follows that
\[
a_2 < \beta\gamma - (\beta + \gamma)^2 = -\beta^2 - \beta\gamma - \gamma^2,
\]
so $a_2$ is negative.

Finally, if one root is equal to 0, then $a_3$ must necessarily
vanish.

By replacing $x$ with $-x$, one concludes that for three
negative roots the necessary and sufficient condition is
\weq{4}{a_1 > 0, \quad a_2 > 0, \quad a_3 > 0.}
We therefore obtain, under the assumption of a positive
discriminant, the following table:
\[
\begin{array}{ccc|c}
a_1 & a_2 & a_3 & \\ \hline
- & + & - & 3\\[2pt]
- & - & - & \\
+ & + & - & \raisebox{1.2em}{$\big\}1$}\\
+ & - & - & \\ \hline
+ & + & + & 0\\[2pt]
- & - & + & \\
- & + & + & \raisebox{1.2em}{$\big\}2$}\\
+ & - & + &
\end{array}
\]
where in the last column the number of positive roots is given.
We can express the result of this consideration as follows:

\gesperrt{For a positive discriminant, the number of
positive roots of the cubic equation is equal to the number of
sign changes in the sequence}
\[
1, \; a_1, \; a_2, \; a_3,
\]
where by a sign change we understand the succession of a
positive and a negative or of a negative and a positive
quantity.
